% Source: https://artofproblemsolving.com/community/c3h2328565_please_help_on_solutionit_does_not_make_sense
% Stable identity: aops-topic-2328565-q01
\textbf{Why $BD=CD$.} The inscribed angles $\angle BAD$ and $\angle DAC$ are equal, and they subtend the chords $BD$ and $DC$ of the same circle. Equal inscribed angles subtend equal chords, so $BD=CD$. ($D$ is the midpoint of the arc $BC$ not containing $A$.)

\begin{asy}
usepackage("amsmath");
size(5.5cm);
defaultpen(fontsize(10pt));
pair B=(0,0), C=(9,0); real ax=11/3; pair A=(ax,sqrt(49-ax^2));
real k=(A.x^2+A.y^2-9*A.x)/(2*A.y); pair O=(4.5,k); real R=abs(O-B);
pair D=(4.5,k-R);
draw(circle(O,R),gray);
draw(A--B--C--cycle,linewidth(0.9)); draw(A--D,blue+linewidth(0.9)); draw(B--D,red); draw(C--D,red);
dot(A); dot(B); dot(C); dot(D);
label("$A$",A,N); label("$B$",B,W); label("$C$",C,E); label("$D$",D,S);
label("$7$",(A+B)/2,NW); label("$8$",(A+C)/2,NE); label("$9$",(B+C)/2+(1.2,0),N);
draw(arc(A,0.9,degrees(B-A),degrees(D-A)),blue); draw(arc(A,1.05,degrees(D-A),degrees(C-A)),blue);
\end{asy}

\textbf{Ptolemy's theorem.} The points $A,B,D,C$ lie on the circle in this order, so for the cyclic quadrilateral $ABDC$ the product of the diagonals equals the sum of the products of opposite sides:
\[AD\cdot BC=AB\cdot CD+AC\cdot BD.\]
With $BD=CD$ this gives $9\cdot AD=7\cdot CD+8\cdot CD=15\cdot CD$, so
\[\dfrac{AD}{CD}=\dfrac{15}{9}=\boxed{\dfrac53}.\]
